% !TEX TS-program = xelatex
% !TEX encoding = UTF-8 Unicode
\documentclass[11pt]{article}
\usepackage[margin=1in]{geometry}

\usepackage{xcolor}
\usepackage{hyperref}
\hypersetup{
  colorlinks=true,
  citecolor=violet
}

\usepackage[citestyle=authoryear, bibstyle=reading, entrykey=false, abstract=false]{biblatex}
\addbibresource{references.bib}

\usepackage{fontspec} 
\setmainfont{Charis SIL}


\title{Thesis Pre-proposal: Allison}
\author{Wren Arbor}
\date{August 15th, 2025}

\begin{document}
\maketitle

If I were to be assigned Allison as my thesis advisor, I would investigate the research question: How does frontness of English /s/ vary within English-speaking trans women's speech with respect to discourse context, audience, and self-identity? Is this behavior better understood through performativity theory or indexicality theory? This is a slightly expanded version of the topic I investigated in my RD qual. The work I did there noticed that in current language and gender scholarship, trans women's average /s/ production as a group has been put on scales of frontness for comparison against average /s/ production of groups of people of other gender/sexuality identities \parencite[see][] {hazenberg_language_2012, podesva_sexuality_2015}, but variation in /s/ production among trans women, both between speakers and within a single speaker's speech, has not been investigated. I reviewed how filling this gap would be illuminating for four primary reasons. 

\begin{enumerate}
\item First, both \textcite{hazenberg_language_2012, podesva_sexuality_2015} find that trans women's /s/ production significantly differs than that of cis people's, but are unable to explain this only by looking at aggregated numerical data-- by taking a closer look at variation of /s/ usage within trans women, a better understanding of this aggregate difference would likely emerge. 

\item Secondly, fronted /s/ has been traditionally regarded as an index of femininity (eg. \cite{campbell-kibler_intersecting_2011, podesva_sexuality_2015, maegaard_s_2015}; also reviewed in \cite{calder_sissy_2019}). Within this model, the fact that (some cis white) gay men use fronted /s/ to index gayness is understood by \textcite{calder_sissy_2019} as a product of indexical orders \parencite[][see]{silverstein_indexical_2003} developing along strata of fractal recursivity \parencite[][see]{kroskrity_language_2000}; gay men become ``men who are women'' through fractal recursivity, and a new indexical order of /s/ emerges at that layer of opposition, going from a link from fronter /s/ to femininity to a link from fronter /s/ in men to gayness. However, trans people cannot be easily fit into this model. If /s/ were simply a feminine and gay-male index, trans people would be expected to pattern just the same as cis people. The fact that they do not (as seen in, eg. \cite{hazenberg_language_2012, podesva_sexuality_2015, zimman_hegemonic_2013}) also cannot be explained as simply another level of fractal recursivity; unlike sexuality identity, which can construe itself ideologically as a layer over and above gender and sex in the fractal recursive tree, this is untenable for transness, which is unavoidably present at all levels. \textcite{calder_whose_2021} points out that this problem reveals the traditional view of /s/ as only feasible within the hegemonic, white, cisgender indexical field, and that an account of how SGS (and race) are produced through language is incomplete without investigations of ``speaking and listening subjects outside of White and cisgender populations'' (43). In this way, investigating trans women's usage of /s/ on their own terms helps to decenter the hegemonic indexical field, and to build a more complete understanding of language's role in constructing SGS and race.

\item Thirdly, by looking outside of the hegemonic indexical field, this research would help to investigate how alternative ways of using variation form. Specifically, it may provide insight into the tension that exists between performativity theory and indexicality theory; if the `meaning' of a variant does not work, or is not straightforwardly taken up by a population, does a new, different meaning necessarily emerge among that population, as the traditional indexicality theory model (specifically the theory of indexical orders) may predict? Or, can the so-called `meaning' be shown to exist as an order-construction within the hegemonic field (cf. \cite{deleuze_november_1987}), and can non-White, non-cisgender people be seen to be using variation not for what it `means', but for what it `does', as would be predicted by performativity theory \parencite[see][]{butler_gender_1990, livia_its_1997}? How can these theories and perspectives be reconciled?

In order to study this, I would likely take a mixed methodology approach. 

As far as data collection, In my qual, I imagined  recruiting a group of transgender women speakers, collecting information on their identities using a pre-survey, and then giving them microphones and recorders, tasking them to record their speech across several contexts, in natural conversation with a diverse set of interlocutors, and to give the interlocutorsn surveys to collect information about their identities. For a thesis, I would supplement this with sociolinguistic interviews with me, focused mostly on collecting explicit metalinguistic commentary on navigating language use while trans. 

As far as analysis, I would perform both quantitative and qualitative analyses. I would first measure COG (maybe also skew, loudness, duration, and other acoustic parameters) for all of the /s/ tokens across produced by my speakers throughout all of my recordings. Then, I would quantitavely analyse the relationship between the identity and interlocutor variables collected for each survey and the /s/ realizations within the associated recording. This would likely be done with a linear mixed effects model. I did this quantitative analysis with fabricated data on my qual. For the qualitative analysis, I would perform discourse analysis on transcriptions of (perhaps a selection of) the recordings and analyze what discourse context can tell us about the choice of an /s/ variant. This was not a part of my qual, but discourse analysis of sociolinguistic interviews is something I'm familiar with from work in Intro Sociolinguistics and Language and Gender.

This is the proposal I'm most excited and confident about, so I really hope I get the chance to do it with Allison! Thank you :)

\printbibliography[title=References with Annotations]


\end{enumerate}
\end{document}